%%%%%%%%%%  scan idx 624  |  folio 611  %%%%%%%%%%
\textbf{Lemme 4.2.} \uline{Sous les conditions de} 4.1, \(f\) \uline{est un morphisme de descente pour la catégorie des faisceaux localement libres~; c'est-à-dire pour tout couple} \(L,M\) \uline{de faisceaux localement libres sur} \(Y\)\uline{, le diagramme}
\[
\begin{tikzcd}[column sep=tiny]
\Hom_Y(L,M) \arrow[r] &
\Hom_X(f^*L,f^*M) \arrow[r,shift left=.7ex] \arrow[r,shift right=.7ex] &
\Hom_{X\times_YX}(g^*L,g^*M)
\end{tikzcd}
\]
\uline{est exact.}

En effet, si \(N=\underline{\Hom}_{{O}_Y}(L,M)\), on doit montrer que la suite
\[
  \Gamma(Y,N)\longrightarrow\Gamma(X,f^*N)
  \rightrightarrows
  \Gamma(X\times_YX,g^*N)
\]
est exacte, ou bien (d'après EGA \(0_1\) 5.4.10.1) que la suite
\[
  \Gamma(Y,N)\longrightarrow\Gamma(Y,f_*{O}_X\otimes_{{O}_Y}N)
  \rightrightarrows
  \Gamma(Y,g_*{O}_{X\times_YX}\otimes_{{O}_Y}N)
\]
est exacte. Mais la suite (*) \({O}_Y\to f_*{O}_X\rightrightarrows g_*{O}_{X\times_YX}\) est exacte par hypothèse~; d'où le lemme, en lui appliquant le foncteur \(\Gamma(N\otimes -)\), qui est exact à gauche, \(N\) étant plat sur \(Y\).

Soit \(L\) un faisceau inversible sur \(X\) muni d'une rigidification \(u\) relativement à \(\alpha\)~; d'après 4.2 on a une équivalence entre la catégorie \(C\) des couples \((M,\varphi)\) d'un faisceau inversible \(M\) sur \(Y\) et d'un isomorphisme \(\varphi:f^*M\xrightarrow{\sim}L\), et la catégorie \(\mathrm{Deff}(S)\) des données de descente effectives sur \(L\) relativement à \(f\). Mais si \(f^*M\simeq L\), comme \(L\) est rigidifié relativement à \(\alpha\), \(M\) est muni canoniquement d'une rigidification relativement à \(\beta\). Par conséquent \(C\) est équivalente à la catégorie discrète définie par l'ensemble des classes de faisceaux inversibles \(\beta\)-rigidifiés \((M,v)\) sur \(Y\), tels que \(f^*(M,v)\simeq(L,u)\). Mais c'est précisément l'ensemble \(Z(S)\) de \(S\)-points de l'image réciproque dans \(\Pic_{Y/S}\) du \(S\)-point de \(\Pic_{X/S}\) défini par \((L,u)\). De plus cette description de \(Z(S)\) est fonctorielle en \(S\). Donc \(Z\) s'identifie au foncteur \(\mathrm{Deff}\), et on peut oublier les rigidifications par la suite.

